\documentclass[11pt]{amsart}
\usepackage{amsmath,amssymb,amsthm,booktabs}
\usepackage[margin=1.15in]{geometry}
\usepackage[hidelinks]{hyperref}
\newtheorem{theorem}{Theorem}[section]
\newtheorem{proposition}[theorem]{Proposition}
\newtheorem{lemma}[theorem]{Lemma}
\newtheorem{corollary}[theorem]{Corollary}
\theoremstyle{remark}
\newtheorem{remark}[theorem]{Remark}
\newcommand{\Z}{\mathbb{Z}}
\newcommand{\R}{\mathbb{R}}
\newcommand{\E}{\mathbb{E}}
\newcommand{\Prob}{\mathbb{P}}
\newcommand{\doi}[1]{\href{https://doi.org/#1}{doi:\nolinkurl{#1}}}
\setlength{\emergencystretch}{2.5em}

\title[Zero speed for excited random walks on the event of non-transience]{Zero speed for excited random walks on the event of non-transience}
\author{Yitzchak Shmalo}
\address{Yitzchak Shmalo, Tel Aviv University and Heights Labs}
\email{yspennstate@gmail.com}
\date{November 2026}
\subjclass[2020]{60K35, 60K37, 60F15}
\keywords{excited random walk, cookie random walk, law of large numbers, self-interacting random walk, zero-one law}

\begin{document}
\begin{abstract}
An excited random walk on $\Z^d$ moves, on its $i$-th visit to a site $x$, according to the $i$-th entry of a stack of
transition probabilities attached to $x$; the stacks are independent and identically distributed over the sites.
Kosygina and Zerner asked whether $X_n/n\to0$ almost surely on the event that $|X_n|$ does not tend to infinity,
assuming only that the stacks are i.i.d.\ and every entry is positive. We prove that it does, in every dimension and with
no further restriction on the stacks, which may be infinite and may push in any direction. We also prove the
excited-walk form of Zerner's law of large numbers for random walk in random environment: if, for a coordinate
direction $e$, the walk tends to $+\infty$ and to $-\infty$ in direction $e$ with probability zero, then $X_n\cdot e/n\to0$
almost surely. With a theorem of the survey this gives the almost sure convergence of $X_n/n$ under the same two
assumptions. The strong Markov property, which Zerner applies at the last visit to a level, is replaced by a sum over the
finitely many paths the walk can follow between its first visit to the level and that visit.
\end{abstract}
\maketitle

\section{Introduction}\label{sec:intro}

Let $E=\{\pm e_1,\dots,\pm e_d\}$ be the unit vectors of $\Z^d$. A \emph{cookie environment} is a family
$\omega=(\omega(x,a,i))_{x\in\Z^d,\,a\in E,\,i\ge1}$ with $\omega(x,a,i)\ge0$ and $\sum_{a\in E}\omega(x,a,i)=1$. The
sequence $\omega(x,\cdot,\cdot)$ is the \emph{stack} at $x$. Given $\omega$ and $y\in\Z^d$, the excited random walk
$(X_n)_{n\ge0}$ starts at $X_0=y$ and moves by
\begin{equation}\label{eq:step}
P_{y,\omega}\bigl(X_{n+1}=X_n+a\,\big|\,X_0,\dots,X_n\bigr)=\omega\bigl(X_n,a,\#\{0\le m\le n:X_m=X_n\}\bigr),\qquad a\in E,
\end{equation}
so on its $i$-th visit to a site the walk uses the $i$-th entry of that site's stack. The environment is random with
law $\Prob$, and $P_y=\E[P_{y,\omega}]$ is the annealed law, where $\E$ is expectation with respect to $\Prob$; $E_y$
is expectation with respect to $P_y$. We use the two assumptions of the survey of Kosygina and
Zerner~\cite{KZ13}:
\begin{itemize}
\item[(IID)] the stacks $\omega(x,\cdot,\cdot)$, $x\in\Z^d$, are independent and identically distributed under $\Prob$;
\item[(EL)] $\Prob$-almost surely, $\omega(x,a,i)>0$ for all $x$, $a$ and $i$.
\end{itemize}
Nothing else is assumed. There is no uniform lower bound in (EL), the entries of one stack need not be independent, and a
stack may have infinitely many entries that differ from the simple random walk step, with drifts of either sign.

For $\ell\in\R^d\setminus\{0\}$ let $A_\ell=\{X_n\cdot\ell\to+\infty\}$, and let
\[
N=\bigl\{|X_n|\not\to\infty\bigr\}
\]
be the event that the walk does not go to infinity, which we call non-transience. Kosygina and Zerner asked the following~\cite[Problem 4.7]{KZ13}.

\medskip\noindent\textbf{Problem.} Assume (IID) and (EL), and let $\ell\in\R^d\setminus\{0\}$. Does $X_n\cdot\ell/n$
converge to $0$ $P_0$-almost surely on $N$?

\medskip
For random walk in an i.i.d.\ elliptic random environment, which is the case in which $\omega(x,a,i)$ does not depend
on $i$, Zerner~\cite[Theorem 1]{Zer02} proved the following: if $e\in E$ and $P_0(A_e\cup A_{-e})=0$, then
$X_n\cdot e/n\to0$ $P_0$-almost surely. In dimension one, Amir, Berger and Orenshtein~\cite[Theorem 1.3]{ABO16} proved
that $X_n/n$ converges almost surely to a deterministic limit for stationary ergodic elliptic stacks on $\Z$, infinite
stacks included; when $P_0(N)>0$ the limit is zero, since on $N$ the walk visits some site infinitely often. We answer
the problem in every dimension.

\begin{theorem}\label{thm:main}
Let $d\ge1$ and assume (IID) and (EL). Then $X_n/n\to0$ $P_0$-almost surely on $N$. If $P_0(N)>0$, then $X_n/n\to0$
$P_0$-almost surely.
\end{theorem}

The main step is the following one-sided statement, which holds for every coordinate direction separately.

\begin{proposition}\label{prop:onesided}
Assume (IID) and (EL), and let $e\in E$. If $P_0(A_e)=0$, then $\limsup_{n}X_n\cdot e/n\le0$ $P_0$-almost surely.
\end{proposition}

Applied to $e$ and to $-e$, Proposition~\ref{prop:onesided} gives the excited-walk form of Zerner's theorem.

\begin{corollary}\label{cor:zerner}
Assume (IID) and (EL), and let $e\in E$. If $P_0(A_e\cup A_{-e})=0$, then $X_n\cdot e/n\to0$ $P_0$-almost surely.
\end{corollary}

When the stacks are constant in $i$ this is~\cite[Theorem 1]{Zer02}. When $P_0(N)>0$, the zero-one
law~\cite[Theorem 3.16]{KZ13} gives $P_0(A_e\cup A_{-e})=0$ for every $e\in E$, and Theorem~\ref{thm:main} follows; the
deduction is in Section~\ref{sec:proof}. With~\cite[Theorem 4.6]{KZ13}, Corollary~\ref{cor:zerner} also gives a law
of large numbers under (IID) and (EL) alone.

\begin{corollary}\label{cor:lln}
Assume (IID) and (EL). Then $X_n/n$ converges $P_0$-almost surely. For $1\le j\le d$, the limit of $X_n\cdot e_j/n$ is
$0$ if $P_0(A_{e_j}\cup A_{-e_j})=0$; otherwise it is a constant $v_{e_j}\ge0$ on $A_{e_j}$ and a constant
$-v_{-e_j}\le0$ on $A_{-e_j}$.
\end{corollary}

\medskip\noindent\textbf{The idea of the proof.} Fix $e\in E$, call $\{x:x\cdot e=m\}$ level $m$, let $T_m$ be the
first time the walk is at level $m$, and for an integer $L\ge1$ let $h_{m,L}$ be the time from $T_m$ to the last visit
to level $m$ before $T_{m+L}$. Zerner's proof has two parts. By Zerner's lemma~\cite[Lemma 3]{Zer02}, a statement about a single path,
if $\limsup_nX_n\cdot e/n>0$ then there is an integer $c$ such that, for every $L$, the upper density
\[
\limsup_{M\to\infty}\frac1{M+1}\,\#\{0\le m\le M:\ h_{m,L}\le c\}
\]
of the levels $m$ with $h_{m,L}\le c$ is bounded below by a positive constant that does not depend on $L$. The second
part bounds the probability that $h_{m,L}\le c$. At its last visit to level $m$ the walk is at one of finitely many
sites $z$, for the $r$-th time, with $r\le c+1$. The $r$-th visit to $z$ is a stopping time, and by the strong Markov
property the walk has law $P_{z,\omega}$ from then on. The probability that it then reaches level $m+L$ without
returning to level $m$ depends only on the environment at levels $\ge m$, which is independent of the path up to time
$T_m$.

For an excited walk the law after the $r$-th visit to $z$ is not $P_{z,\omega}$. It depends on which entries of the
stacks at $z$, and at the other sites visited since $T_m$, have been used, and this is not determined by $z$ and $r$.
It is determined by the path of the walk from $T_m$ to that visit. Let $\Pi_c$ be the finite set of nearest-neighbour
paths $p=(p_0,\dots,p_s)$ with $p_0=0$, $s\le c$ and $p_s\cdot e=0$, and let $\omega^p$ be the environment obtained
from $\omega$ by removing from each stack the entries used along $p$. Bounding by one the probability that the walk
follows $p$ after $T_m$, we obtain under (IID) alone, for all $m\ge0$ and $L>c$, on $\{T_m<\infty\}$,
\[
P_0\bigl(T_{m+L}<\infty,\ h_{m,L}\le c\,\big|\,X_0,\dots,X_{T_m}\bigr)\le
\sum_{p\in\Pi_c}\E\Bigl[P_{p_s,\omega^p}\bigl(T_L<\infty,\ X_n\cdot e>0\text{ for }1\le n\le T_L\bigr)\Bigr]
\]
(Lemma~\ref{lem:fresh}). The right side depends neither on $m$ nor on the path up to time $T_m$. Each term decreases,
as $L\to\infty$, to the expected probability that the walk in $\omega^p$ started at $p_s$ stays above level $0$ and
reaches every level, and under (EL) this is zero when $P_0(A_e)=0$ (Lemma~\ref{lem:prefix}). Hence, almost surely, the
upper density of the levels $m$ with $h_{m,L}\le c$ is at most the right side, which tends to $0$ as $L\to\infty$ for
each $c$, and Zerner's lemma gives $\limsup_nX_n\cdot e/n\le0$.

\medskip\noindent\textbf{Related work.} Excited random walk was introduced by Benjamini and Wilson~\cite{BW03}, with one
cookie of drift in direction $e_1$ at each site. Zerner~\cite{Zer06} and Menshikov, Popov, Ram\'irez and
Vachkovskaia~\cite{MPRV12} studied walks whose drifts point into a half-space, and B\'erard and Ram\'irez~\cite{BR07} proved
a central limit theorem for the Benjamini--Wilson walk; for $d\ge2$ the walk in these models is transient, and $P_0(N)=0$. The laws of
large numbers of Holmes~\cite{Hol12} and Pham~\cite{Pha15} in high dimensions assume simple transverse steps.
Theorem~4.8 of the survey~\cite{KZ13} gives $\lim_nX_n/n\in\{v,-cv\}$ for some $v\in\R^d$ and $c\ge0$ when
$P_0(A_{\ell_i}\cup A_{-\ell_i})=1$ for every vector of a basis $\ell_1,\dots,\ell_d$ of $\R^d$; the survey attributes
this assumption to Problem~4.7. By the zero-one law~\cite[Theorem 3.16]{KZ13} each of these probabilities is $0$ or
$1$, and Corollary~\ref{cor:zerner} gives $X_n\cdot e/n\to0$ for a coordinate direction $e$ in which it is $0$. In dimension one, besides~\cite{ABO16}, infinite stacks
were treated by Kosygina and Peterson~\cite{KP17} for Markovian stacks.

A table of notation is in Appendix~\ref{app:notation}.

\section{Preliminaries}\label{sec:prelim}

\subsection{Results from the survey}
We use three results stated in~\cite{KZ13}: (R1) and (R3) under (IID) and (EL), and (R2) under (IID) and
$\E[\omega(0,e,1)]>0$ for all $e\in E$, which (EL) implies. The survey proves (R1) and (R2) and sketches the proof of
(R3), which follows that of Zerner and Merkl~\cite[Proposition 3]{ZM01} for random walk in random environment.

\begin{itemize}
\item[(R1)] \cite[Theorem 2.5]{KZ13} The range $\{X_n:n\ge0\}$ is either $P_0$-almost surely finite or
$P_0$-almost surely infinite.
\item[(R2)] \cite[Lemma 3.18]{KZ13} If the range is $P_0$-almost surely infinite, then for every
$\ell\in\R^d\setminus\{0\}$ and all real $u<w$, the walk almost surely does not stay in the half-space
$\{x:x\cdot\ell\le w\}$ forever while visiting the slab $\{x:u\le x\cdot\ell\le w\}$ infinitely often.
\item[(R3)] \cite[Theorem 3.16]{KZ13} For every $\ell\in\R^d\setminus\{0\}$, $P_0(A_\ell\cup A_{-\ell})\in\{0,1\}$.
\end{itemize}
We use (R2) only with $\ell$ a coordinate vector and $u,w$ integers.

\begin{lemma}\label{lem:onesidedpaths}
Assume (IID) and (EL), suppose the range is $P_0$-almost surely infinite, and let $e\in E$. Then
\[
P_0\Bigl(\inf_n X_n\cdot e>-\infty,\ X_n\cdot e\not\to+\infty\Bigr)=0 .
\]
\end{lemma}

\begin{proof}
On the event in question there are integers $a<b$ such that $X_n\cdot e\ge a$ for all $n$ and $X_n\cdot e\le b$ for
infinitely many $n$. For fixed $a<b$ this is the event of (R2) with $\ell=-e$, $w=-a$ and $u=-b$, which is null. Take the
union over the countably many pairs $(a,b)$.
\end{proof}

\subsection{Zerner's lemma}
The quantities $T_m$ and $h_{m,L}$ of Section~\ref{sec:intro} are defined for every nearest-neighbour path
$(x_n)_{n\ge0}$ in $\Z^d$ with $x_0\cdot e=0$, a direction $e\in E$ and integers $m\ge0$, $L\ge1$; in formulas,
\[
T_m=\inf\{n\ge0:x_n\cdot e=m\},\qquad I_{m,L}=\{n:T_m\le n<T_{m+L},\ x_n\cdot e=m\},
\]
and $h_{m,L}=\max I_{m,L}-T_m$ when $T_{m+L}<\infty$, with the convention $h_{m,L}=\infty$ otherwise. For an integer
$c\ge0$ put
\[
F_{M,L}(c)=\frac1{M+1}\sum_{m=0}^{M}\mathbf 1\{h_{m,L}\le c\}.
\]
Since $h_{m,L}$ does not decrease in $L$, $F_{M,L}(c)$ does not increase in $L$. The next lemma is Zerner's; it
concerns a single path, and its proof in~\cite{Zer02} uses no property of the environment.

\begin{lemma}[{\cite[Lemma 3]{Zer02}}]\label{lem:counting}
Let $e\in E$ and let $(x_n)_{n\ge0}$ be a nearest-neighbour path in $\Z^d$ with $x_0\cdot e=0$. If
$\limsup_n x_n\cdot e/n>0$, then there is an integer $c\ge0$ with
\[
\inf_{L\ge1}\,\limsup_{M\to\infty}F_{M,L}(c)>0 .
\]
\end{lemma}

\section{The estimate at the first visit to a level}\label{sec:fresh}

Fix $e\in E$ and recall $\Pi_c$ and $\omega^p$ from Section~\ref{sec:intro}; in formulas, for $p\in\Pi_c$ let
$k_p(x)=\#\{0\le t<s:p_t=x\}$ be the number of departures from $x$ along $p$, and $\omega^p(x,a,i)=\omega(x,a,i+k_p(x))$.
The
endpoint $p_s$ has not been left along $p$ after its last visit, so its next entry is not removed. For $L>c$ define
\[
q_L^p(\omega)=P_{p_s,\omega^p}\bigl(T_L<D\bigr),\qquad D=\inf\{n\ge1:X_n\cdot e\le0\},
\]
for the walk started at $p_s$. On $\{T_L<D\}$ the walk leaves level $0$ at
its first step and stays strictly above it until it reaches level $L$, so $q_L^p$ depends only on the stacks at sites
$x$ with $x\cdot e\ge0$, even if $p$ itself visits negative levels. As $L$ increases, $q_L^p$ decreases to
\[
q_\infty^p(\omega)=P_{p_s,\omega^p}\Bigl(X_n\cdot e>0\text{ for all }n\ge1\text{ and }\sup_nX_n\cdot e=\infty\Bigr),
\]
the probability of reaching every level without returning to level $0$. Put
\[
a_L(c)=\sum_{p\in\Pi_c}\E\bigl[q_L^p\bigr],\qquad L>c .
\]

\begin{lemma}\label{lem:prefix}
Assume (IID) and (EL), and suppose $P_0(A_e)=0$. Then for every $c\ge0$, $a_L(c)\to0$ as $L\to\infty$.
\end{lemma}

\begin{proof}
Fix $p\in\Pi_c$ and let $w_p(\omega)=P_{0,\omega}\bigl((X_0,\dots,X_s)=p\bigr)$, a finite product of entries of
$\omega$, which is positive $\Prob$-almost surely by (EL). After following $p$ the walk sits at $p_s$, and the stacks it has
not yet used are those of $\omega^p$, so by~\eqref{eq:step}, $w_p\,q_\infty^p$ is the probability under $P_{0,\omega}$ that the walk follows $p$
and afterwards stays strictly above level $0$ and reaches every level. By (R1) the range is $P_0$-almost surely finite or
$P_0$-almost surely infinite. In the first case this event is $P_0$-null. In the second, since $X_n\cdot e\ge-c$ for all $n$ on this event, Lemma~\ref{lem:onesidedpaths} shows
that it is contained in $A_e$ up to a $P_0$-null set. In both cases $\E[w_pq_\infty^p]\le P_0(A_e)=0$, and since
$w_p>0$ almost surely, $q_\infty^p=0$ almost surely. By dominated convergence $\E[q_L^p]\to0$, and $\Pi_c$ is finite.
\end{proof}

Let $\mathcal F_n=\sigma(X_0,\dots,X_n)$; for a stopping time $T$, $\mathcal F_T$ is the $\sigma$-field of events
determined by the path up to time $T$. For $m\ge0$, $c\ge0$ and $L>c$ let
\[
B_{m,L,c}=\{T_{m+L}<\infty,\ h_{m,L}\le c\}.
\]

\begin{lemma}\label{lem:fresh}
Assume (IID). Let $m\ge0$, $c\ge0$ and $L>c$. Then $P_0(B_{m,L,c}\mid\mathcal F_{T_m})\le a_L(c)$ on $\{T_m<\infty\}$.
\end{lemma}

\begin{proof}
Fix a finite path $\eta=(\eta_0,\dots,\eta_t)$ with $\eta_0=0$ whose first visit to level $m$ is at step $t$, let
$x=\eta_t$, and let $w_\eta(\omega)=P_{0,\omega}\bigl((X_0,\dots,X_t)=\eta\bigr)$. Before time $t$ the walk leaves only
sites at levels below $m$, so $w_\eta$ is a product of entries of stacks at sites $y$ with $y\cdot e<m$.

On $B_{m,L,c}\cap\{(X_0,\dots,X_t)=\eta\}$ let $t+s$ be the last visit to level $m$ before $T_{m+L}$, so that
$s=h_{m,L}\le c$, and let $p=(X_{t+r}-x)_{0\le r\le s}$; then $p\in\Pi_c$. Conversely, if the walk follows $\eta$, then
$x+p$ for some $p\in\Pi_c$, and then reaches level $m+L$ before returning to level $m$, it lies in $B_{m,L,c}$ with
$h_{m,L}=s$, because the path $x+p$ has at most $c<L$ steps and does not reach level $m+L$. These events are disjoint
for different $p$: if $p$ is a proper initial piece of $p'$, a walk that follows $x+p'$ is at level $m$ at a step
after $t+s$ and before $T_{m+L}$.

Let $\omega'$ be the environment with the entries used by $\eta$ removed, and let
$(\theta_x\omega')(z,a,i)=\omega'(x+z,a,i)$. After following $\eta$ and then $x+p$ the walk is at $x+p_s$, and
by~\eqref{eq:step} it then reaches level $m+L$ before returning to level $m$ with probability
$q_L^p(\theta_x\omega')$. As noted after the definition of $q_L^p$, it reads only stacks at levels $\ge m$ after the shift, and $\eta$
uses no entry of these stacks, so $q_L^p(\theta_x\omega')=q_L^p(\theta_x\omega)$. The probability of following $\eta$
and then $x+p$ is at most $w_\eta(\omega)$, so
\begin{align*}
P_{0,\omega}\bigl(B_{m,L,c},\ (X_0,\dots,X_t)=\eta\bigr)
&=\sum_{p\in\Pi_c}P_{0,\omega}\bigl((X_0,\dots,X_t)=\eta,\ (X_{t+r}-x)_{0\le r\le s}=p\bigr)\,q_L^p(\theta_x\omega)\\
&\le w_\eta(\omega)\sum_{p\in\Pi_c}q_L^p(\theta_x\omega),
\end{align*}
where $s$ is the number of steps of $p$.
The factor $w_\eta$ depends only on the stacks at levels $<m$ and the sum only on the stacks at levels $\ge m$, which
are independent by (IID); by translation invariance $\E[q_L^p\circ\theta_x]=\E[q_L^p]$. Taking expectations,
$P_0\bigl(B_{m,L,c},\ (X_0,\dots,X_t)=\eta\bigr)\le a_L(c)\,P_0\bigl((X_0,\dots,X_t)=\eta\bigr)$. An event of
$\mathcal F_{T_m}$ contained in $\{T_m<\infty\}$ is a disjoint union of events $\{(X_0,\dots,X_t)=\eta\}$ of this
kind, and summing over them gives the claim.
\end{proof}

\section{Proof of Proposition~\ref{prop:onesided}}\label{sec:onesided}

Let $P_0(A_e)=0$.

Fix $c\ge0$, $L>c$ and a residue $r\in\{0,\dots,L-1\}$, and put $m_k=r+kL$, $\mathcal G_k=\mathcal F_{T_{m_k}}$ and
$Z_k=\mathbf 1_{B_{m_k,L,c}}$. The event
$B_{m_k,L,c}$ is determined by the path up to $T_{m_k+L}=T_{m_{k+1}}$, so $Z_k$ is $\mathcal G_{k+1}$-measurable, and
$Z_k-E_0[Z_k\mid\mathcal G_k]$ is a martingale difference sequence bounded by $1$. By the Azuma--Hoeffding inequality,
$P_0\bigl(|\sum_{k<K}(Z_k-E_0[Z_k\mid\mathcal G_k])|>\varepsilon K\bigr)\le2e^{-\varepsilon^2K/2}$, so by the
Borel--Cantelli lemma $K^{-1}\sum_{k<K}(Z_k-E_0[Z_k\mid\mathcal G_k])\to0$ $P_0$-almost surely. With
Lemma~\ref{lem:fresh},
\[
\limsup_{K\to\infty}\frac1K\sum_{k<K}Z_k\le a_L(c)\qquad P_0\text{-a.s.}
\]
If some $T_{m_k}$ is infinite then $Z_j=0$ for all $j\ge k$ and the bound holds trivially. Averaging over the $L$
residues and intersecting over the countably many pairs $(c,L)$, almost surely, for all $c\ge0$ and $L>c$,
\begin{equation}\label{eq:freq}
\limsup_{M\to\infty}\frac1{M+1}\sum_{m=0}^{M}\mathbf 1_{B_{m,L,c}}\le a_L(c).
\end{equation}

Now take a path in the almost sure event of~\eqref{eq:freq} with $\limsup_nX_n\cdot e/n>0$. Since $h_{m,L}=\infty$ when
$T_{m+L}=\infty$, $B_{m,L,c}=\{h_{m,L}\le c\}$, and the left side of~\eqref{eq:freq} is $\limsup_MF_{M,L}(c)$. By
Lemma~\ref{lem:prefix}, $\inf_{L>c}\limsup_MF_{M,L}(c)=0$ for every $c$; since $F_{M,L}(c)$ does not increase in $L$,
the infimum over all $L\ge1$ is also zero. This contradicts Lemma~\ref{lem:counting}. Hence
$\limsup_nX_n\cdot e/n\le0$ almost surely. \qed

\section{Proofs of Corollaries~\ref{cor:zerner} and~\ref{cor:lln} and Theorem~\ref{thm:main}}\label{sec:proof}

\begin{proof}[Proof of Corollary~\ref{cor:zerner}]
If $P_0(A_e\cup A_{-e})=0$, then $P_0(A_e)=P_0(A_{-e})=0$. Proposition~\ref{prop:onesided}, applied to $e$ and to $-e$,
gives $\limsup_nX_n\cdot e/n\le0$ and $\liminf_nX_n\cdot e/n\ge0$ almost surely.
\end{proof}

\begin{proof}[Proof of Theorem~\ref{thm:main}]
If $P_0(N)=0$ there is nothing to prove. Suppose $P_0(N)>0$ and let $j\in\{1,\dots,d\}$. On $A_{e_j}\cup A_{-e_j}$ we
have $|X_n|\to\infty$, so $N$ is disjoint from it and $P_0(A_{e_j}\cup A_{-e_j})<1$. By (R3) this probability is $0$,
and Corollary~\ref{cor:zerner} gives $X_n\cdot e_j/n\to0$ almost surely. Hence $X_n/n\to0$ $P_0$-almost surely, and in
particular on $N$.
\end{proof}

\begin{proof}[Proof of Corollary~\ref{cor:lln}]
Let $1\le j\le d$. By (R3), $P_0(A_{e_j}\cup A_{-e_j})\in\{0,1\}$. If it is $0$, Corollary~\ref{cor:zerner} gives
$X_n\cdot e_j/n\to0$. If it is $1$, apply~\cite[Theorem 4.6]{KZ13} to each of $e_j$ and $-e_j$ for which
$A_{\pm e_j}$ has positive probability: on $A_{\ell}$, $X_n\cdot\ell/n$ converges to a constant $v_\ell\ge0$. Since
$A_{e_j}\cup A_{-e_j}$ has probability $1$, $X_n\cdot e_j/n$ converges $P_0$-almost surely in this case too. This holds
for every $j$, so $X_n/n$ converges $P_0$-almost surely.
\end{proof}

\begin{remark}\label{rem:trap}
The event $N$ need not be the event that every site is visited infinitely
often. For instance, in $d=2$ give each site an independent uniformly chosen $a\in E$, and let the $i$-th entry of its stack put
mass $1-3\varepsilon_i$ on $a$ and $\varepsilon_i=2^{-i-4}$ on each other direction. With positive probability the
directions at $0$ and $e_1$ point at each other and the walk moves back and forth between these two sites forever, since
\[
\frac1{16}\prod_{i\ge1}\bigl(1-3\cdot2^{-i-4}\bigr)^2>0 .
\]
This environment satisfies (IID) and (EL). Its range is finite with positive probability, hence $P_0$-almost surely
finite by (R1), and $P_0(N)=1$.
\end{remark}

\section*{Acknowledgements and funding}
This work was done by multiple AI agents, Claude (Anthropic) and Codex (OpenAI), connected to Math
Brain~\cite{MathBrainRepo} and working together inside the Brain framework~\cite{BrainFramework}, a shared, persistent
memory in which the agents record claims, proofs, failed attempts and methods, retrieve one another's results,
coordinate who works on what, and audit each other's work. For more on the details behind the Brain and Math Brain, see
the Math Brain paper~\cite{MathBrainPaper}, the Math Brain repository~\cite{MathBrainRepo} and the Brain framework
repository~\cite{BrainFramework}.

The research presented in this paper was supported by the European Research Council (ERC) under the European Union's
Horizon Europe research and innovation programme (grant agreement No.~101041711), by the Simons Foundation as part of
the Collaboration on the Mathematical and Scientific Foundations of Deep Learning, by Heights Labs, by Convex Nexus
Capital, by the Israel Science Foundation (grant number 2258/19), and by the Israel Science Foundation (ISF Grant
4101/25). Computational resources were provided in part by Google Cloud through its education credits program (credit
530550103).

\section*{Declaration of competing interest}
The author holds equity in Heights Labs and in Convex Nexus Capital and has a management role at Convex Nexus Capital.
These entities are included in the funding declaration above.

\section*{Declaration of generative AI and AI-assisted technologies}
A Codex agent (OpenAI) proved Theorem~\ref{thm:main}; a second Codex agent checked the proof as a referee; a Claude
agent (Anthropic) checked the steps and prepared the manuscript. The author directed the work and is
responsible for the content of the paper and for decisions concerning its dissemination.

\section*{Code availability}
Exact finite-horizon checks of the path sum in the proof of Lemma~\ref{lem:fresh}, checks of $F_{M,L}(c)$ on
random paths, and a simulation of the example in Remark~\ref{rem:trap} are in the repository
\begin{center}\url{https://github.com/yspennstate/excited-walk-zero-speed-non-transience}.\end{center}

\begin{thebibliography}{99}
\bibitem{ABO16} G.~Amir, N.~Berger and T.~Orenshtein, Zero-one law for directional transience of one dimensional excited
random walks, \emph{Ann. Inst. Henri Poincar\'e Probab. Stat.} \textbf{52} (2016), 47--57, \doi{10.1214/14-AIHP615}.
\bibitem{BR07} J.~B\'erard and A.~F.~Ram\'irez, Central limit theorem for the excited random walk in dimension
$d\ge2$, \emph{Electron. Commun. Probab.} \textbf{12} (2007), 303--314, \doi{10.1214/ECP.v12-1317}.
\bibitem{BW03} I.~Benjamini and D.~B.~Wilson, Excited random walk, \emph{Electron. Commun. Probab.} \textbf{8} (2003),
86--92, \doi{10.1214/ECP.v8-1072}.
\bibitem{Hol12} M.~Holmes, Excited against the tide: a random walk with competing drifts, \emph{Ann. Inst. Henri
Poincar\'e Probab. Stat.} \textbf{48} (2012), 745--773, \doi{10.1214/11-AIHP434}.
\bibitem{KP17} E.~Kosygina and J.~Peterson, Excited random walks with Markovian cookie stacks, \emph{Ann. Inst. Henri
Poincar\'e Probab. Stat.} \textbf{53} (2017), 1458--1497, \doi{10.1214/16-AIHP761}.
\bibitem{KZ13} E.~Kosygina and M.~P.~W.~Zerner, Excited random walks: results, methods, open problems, \emph{Bull. Inst.
Math. Acad. Sin. (N.S.)} \textbf{8} (2013), 105--157, arXiv:1204.1895.
\bibitem{MPRV12} M.~Menshikov, S.~Popov, A.~F.~Ram\'irez and M.~Vachkovskaia, On a general many-dimensional excited
random walk, \emph{Ann. Probab.} \textbf{40} (2012), 2106--2130, \doi{10.1214/11-AOP678}.
\bibitem{Pha15} C.-D.~Pham, Monotonicity and regularity of the speed for excited random walks in higher dimensions,
\emph{Electron. J. Probab.} \textbf{20} (2015), paper no.~72, \doi{10.1214/EJP.v20-3522}.
\bibitem{Zer02} M.~P.~W.~Zerner, A non-ballistic law of large numbers for random walks in i.i.d.\ random environment,
\emph{Electron. Commun. Probab.} \textbf{7} (2002), 191--197, \doi{10.1214/ECP.v7-1060}.
\bibitem{Zer06} M.~P.~W.~Zerner, Recurrence and transience of excited random walks on $\Z^d$ and strips,
\emph{Electron. Commun. Probab.} \textbf{11} (2006), 118--128, \doi{10.1214/ECP.v11-1200}.
\bibitem{ZM01} M.~P.~W.~Zerner and F.~Merkl, A zero-one law for planar random walks in random environment,
\emph{Ann. Probab.} \textbf{29} (2001), 1716--1732, \doi{10.1214/aop/1015345769}.
\bibitem{MathBrainRepo} Y.~Shmalo, \emph{Math Brain: public snapshot, code and tests}, 2026,
\url{https://github.com/yspennstate/math-brain-public}.
\bibitem{MathBrainPaper} Y.~Shmalo, \emph{Math Brain: shared, auditable memory for multi-agent mathematical research},
October 2026, \url{https://github.com/yspennstate/math-brain-public/tree/main/papers/math-brain-2026-10-24}.
\bibitem{BrainFramework} Y.~Shmalo, \emph{Brain Framework}, public software repository, 2026,
\url{https://github.com/yspennstate/brain-framework-public}.
\end{thebibliography}

\appendix
\section{Notation}\label{app:notation}
\begin{center}
\small
\begin{tabular}{@{}lp{0.6\textwidth}l@{}}
\toprule
Symbol & Meaning & First appears\\
\midrule
$E$ & the $2d$ unit vectors of $\Z^d$ & Section~\ref{sec:intro}\\
$\omega(x,a,i)$ & probability of step $a$ on the $i$-th visit to $x$ (an entry of the stack at $x$) & Section~\ref{sec:intro}\\
$X_n$ & the excited random walk; its range is $\{X_n:n\ge0\}$ & Section~\ref{sec:intro}\\
$P_{y,\omega}$, $P_y$, $\Prob$ & quenched law, annealed law, law of the environment & Section~\ref{sec:intro}\\
(IID), (EL) & the two assumptions on $\Prob$ & Section~\ref{sec:intro}\\
$\E$, $E_y$ & expectation with respect to $\Prob$ and to $P_y$ & Section~\ref{sec:intro}\\
$A_\ell$ & the event $X_n\cdot\ell\to+\infty$ & Section~\ref{sec:intro}\\
$N$ & the event $|X_n|\not\to\infty$ & Section~\ref{sec:intro}\\
level $m$ & the set $\{x:x\cdot e=m\}$ & Section~\ref{sec:intro}\\
$T_m$, $h_{m,L}$ & first time at level $m$; time from $T_m$ to the last visit to level $m$ before $T_{m+L}$ & Section~\ref{sec:intro}\\
$\Pi_c$ & nearest-neighbour paths from $0$ with at most $c$ steps that end at level $0$ & Section~\ref{sec:intro}\\
$\omega^p$ & the environment with the entries used along $p$ removed & Section~\ref{sec:intro}\\
$I_{m,L}$ & visits to level $m$ before level $m+L$ & Section~\ref{sec:prelim}\\
$F_{M,L}(c)$ & fraction of levels $m\le M$ with $h_{m,L}\le c$ & Section~\ref{sec:prelim}\\
$k_p$ & departure counts of $p$ & Section~\ref{sec:fresh}\\
$q_L^p$, $q_\infty^p$, $D$ & probability of reaching level $L$ before time $D$; its limit; the first time $n\ge1$ with $X_n\cdot e\le0$ & Section~\ref{sec:fresh}\\
$a_L(c)$ & $\sum_{p\in\Pi_c}\E[q_L^p]$ & Section~\ref{sec:fresh}\\
$w_p$ & quenched probability of following $p$ & Lemma~\ref{lem:prefix}\\
$\mathcal F_n$, $\mathcal F_T$ & the natural filtration; the $\sigma$-field at a stopping time & Section~\ref{sec:fresh}\\
$B_{m,L,c}$ & the event $T_{m+L}<\infty$, $h_{m,L}\le c$ & Section~\ref{sec:fresh}\\
$\eta$, $w_\eta$, $\omega'$, $\theta_x$ & a path up to its first visit to level $m$; quenched probability of following $\eta$; the environment with the entries used by $\eta$ removed; the shift by $x$ & Lemma~\ref{lem:fresh}\\
$m_k$, $\mathcal G_k$, $Z_k$ & the levels $r+kL$; $\mathcal F_{T_{m_k}}$; the indicator of $B_{m_k,L,c}$ & Section~\ref{sec:onesided}\\
\bottomrule
\end{tabular}
\end{center}

\end{document}
